\chapter{Software Tools Used}

Five tools did real work in this project rather than merely hosting it. For each,
this chapter states the purpose, the specific use in \project{}, and the artefact
that use produced. Screenshots are referenced through the image registry, so a
replacement capture only needs a path change in \code{images.tex}.

\section{Git and GitHub}

\textbf{Purpose.} Distributed version control and a remote home for the repository,
including history, branching and review.

\textbf{Use in \project{}.} The repository holds all three deliverables so that a
change to the API and the client that consumes it can be reviewed together. Work
was done on feature branches and merged into the integration branch
(\code{finish-backend}); the history records the progression from backend scaffolding
to the complete system. Commit messages state the reason for a change, which makes
the history usable as a design log. Git also enforces hygiene: \code{.env},
\code{target/} and the generated \code{reports/} tree are ignored, so no secret or
build artefact enters version control.

\reportimage[0.92]{\ImgToolGit}{Git history and branch state in the repository.}{fig:tool-git}

\section{Docker and Docker Compose}

\textbf{Purpose.} Reproducible build and runtime environments, and a single command
that brings up the database, the application and the test runner.

\textbf{Use in \project{}.} The backend is built in a multi-stage image: a Maven
stage resolves dependencies and packages \code{target/app.jar}, a runtime stage runs
it as a non-root user, and a third stage reuses the build image as the test runner.
Compose defines PostgreSQL with a health check, the backend with a health check and a
dependency on the database, and a test-only profile that runs \code{mvn verify} and
writes reports back to the host. The test project uses its own compose project name,
ports and volume, so a verification run never disturbs a development stack.

\reportimage[0.92]{\ImgToolDocker}{Compose stack in operation and the test workflow output.}{fig:tool-docker}

\section{Postman}

\textbf{Purpose.} Request-level API exploration and a durable, executable record of
the API contract.

\textbf{Use in \project{}.} The collection stored at
\path{src/test/postman/dwara-auth.postman_collection.json} concentrates on the
authentication controller: registration, duplicate registration, login, a wrong
password, the authenticated profile call and the unauthenticated profile call. It
contains six requests and nineteen assertions, and it chains state: the registration
request generates a unique e-mail and stores the returned user id, the login request
stores the token, and the profile request proves that the token resolves to the same
account. A collection pre-request script makes the suite repeatable, so the same
collection can be run from the Postman runner or from Newman without leaving
duplicate accounts behind.

\reportimage[0.92]{\ImgToolPostman}{The authentication collection with its assertions in the Postman runner.}{fig:tool-postman}

\section{Visual Studio Code}

\textbf{Purpose.} One editor for Java, TypeScript, shell and Arduino sources, with
language support that catches errors before a build.

\textbf{Use in \project{}.} VS Code hosted the backend, the Expo application and the
firmware side by side. The Java extension provided navigation across controller,
service and repository layers; the TypeScript server enforced the types shared
between the API client and the screens; the integrated terminal ran the Maven
wrapper, the Compose stack and \code{arduino-cli}. Keeping all three surfaces in one
workspace made it practical to change a DTO and its TypeScript type in the same
edit.

\reportimage[0.92]{\ImgToolVsCode}{The monorepo open in VS Code, with the backend and the application in one
  workspace.}{fig:tool-vscode}

\section{Arduino CLI and Arduino IDE}

\textbf{Purpose.} Compile and flash the ESP32 firmware, and read the serial monitor
during bench testing.

\textbf{Use in \project{}.} The firmware is compiled headlessly with
\code{arduino-cli} against the ESP32 core and the Huge APP partition scheme, which
makes a compile check possible without opening a GUI:

\begin{lstlisting}[style=bash,caption={Firmware compile check used after every firmware
change.},label={lst:compile}]
arduino-cli compile --fqbn esp32:esp32:esp32:PartitionScheme=huge_app \
  firmware/dwara-gate
\end{lstlisting}

The IDE was used for flashing and for the serial monitor at 115200 baud, where the
boot banner reports the selected transport, the servo pin and the advertised BLE
name. The same monitor carries the framed serial-bridge traffic used in the Wi-Fi-less
demonstration.

\reportimage[0.92]{\ImgToolArduino}{Firmware compile and the gate's serial output during a bench run.}{fig:tool-arduino}

\section{Summary}

\begin{table}[htbp]
  \centering
  \small
  \begin{tabularx}{\textwidth}{@{}lXXX@{}}
    \toprule
    \textbf{Tool} & \textbf{Purpose} & \textbf{Phase of SDLC} &
    \textbf{Usage in this project} \\
    \midrule
    Git / GitHub & Version control and collaboration & All phases &
      Monorepo, feature branches, reviewable history, ignored secrets and build
      output \\
    Docker / Compose & Reproducible environments & Implementation, testing,
      deployment & Multi-stage backend image, PostgreSQL service, test profile,
      bind-mounted reports \\
    Postman & API contract testing & Testing & Six-request authentication
      collection with nineteen chained assertions; runnable in the app or Newman \\
    VS Code & Editing and static analysis & Design, implementation, testing &
      Java navigation, TypeScript checking, integrated terminal for Maven, Compose
      and Arduino CLI \\
    Arduino CLI / IDE & Firmware build and bench diagnostics & Implementation,
      testing & Headless compile check, flashing, serial monitor and serial-bridge
      traffic \\
    \bottomrule
  \end{tabularx}
  \caption{Tool summary.}
  \label{tab:tool-summary}
\end{table}
